%%%%%%%%%%  scan idx 229  |  folio 216  %%%%%%%%%%
\[
  \underline C'(S,C^\infty(V,F^\bullet))
  \longrightarrow \underline C'(S,C_2^r(V,F^\bullet))
\]
est un
quasi-isomorphisme de complexes de préfaisceaux, i.e.\ induit des isomorphismes sur les préfaisceaux de cohomologie.

(b) Pour $r\in\Z$, le complexe $\underline F_r^\bullet$ est parfait.

\emph{Preuve de} (b). L'ellipticité de $E^\bullet$ implique ([16]~XI Th.6\footnote{qui se généralise aisément aux complexes elliptiques.}) que, pour tout $s\in S$, la différentielle de $(F_r^\bullet)_s$ est d'image fermée et que les $H^i((F_r^\bullet)_s)$ sont des espaces vectoriels de dimension finie. Autrement dit, pour employer la terminologie de [11], $F_r^\bullet$ est un complexe quasi-acyclique direct de fibrés hilbertiens. Par suite, d'après ([12]~2.3), $F_r^\bullet$ est localement homotopiquement équivalent à un complexe borné de fibrés vectoriels de rang fini. Comme le foncteur «~faisceau de sections~» est additif (3.0.5), $\underline F_r^\bullet$ est localement homotopiquement équivalent à un complexe borné de faisceaux localement libres de type fini, donc est parfait.

Pour la preuve de (a), nous aurons besoin de l'assertion intermédiaire suivante~:

(a') Pour $p\geq q$, le morphisme canonique
\[
  \underline F_p^\bullet\longrightarrow \underline F_q^\bullet
\]
est un quasi-isomorphisme.

\emph{Preuve de} (a'). Notons $M_{pq}^\bullet$ le complexe de fibrés hilbertiens cône du morphisme $F_p^\bullet\longrightarrow F_q^\bullet$. L'ellipticité de $E$ implique que, pour tout $s\in S$, le morphisme $(F_p^\bullet)_s\longrightarrow (F_q^\bullet)_s$ est un quasi-isomorphisme ([16]~XI Cor.~2 of Th.~6). En d'autre termes, $M_{pq}^\bullet$ est un complexe acyclique direct. Il en résulte ([11]~2.3) que $M_{pq}^\bullet$ est localement homotopiquement trivial, donc que $\underline M_{pq}^\bullet$, qui n'est autre que le cône de $\underline F_p^\bullet\longrightarrow\underline F_q^\bullet$, est acyclique, cqfd.

\emph{Preuve de} (a). Il s'agit de montrer que, pour tout ouvert relativement compact $U$ de $S$, le morphisme
